\PuzzleChapterIntro{The Oath of Two Echoes}{Algebra \textbullet\ Hidden Symmetries \textbullet\ Unique Integer}{This puzzle relies on the relationship between the coefficients of a polynomial and the symmetric sums of its roots. You must construct a linear system to solve for the unknown products without finding the roots explicitly.}
\Lore{
The Algebra Vault sells “bonds” on roots the way a bazaar sells grain: by weight, not by story.

A quartic \(f(x)\in\mathbb Z[x]\) is posted on the wall, and four roots \(r_1,r_2,r_3,r_4\) are recorded—two rational, two irrational, all distinct.
The clerk will not sell the roots themselves. The clerk sells only two products:
\(B_1=r_1r_2\) and \(B_2=r_3r_4\).

The price is old-fashioned: if \(B=p/q\) in lowest terms, you pay \(p+q\) coins.
}

\Problem
The Echo Gauge is governed by the plate
\[
169x^4 - 39x^3 + 428x^2 + 89x + 203.
\]
It has exactly four echoes \(r_1,r_2,r_3,r_4\) (not necessarily real): these are the four values of \(x\) that make the gauge’s hum vanish,
i.e. the four roots of the equation
\[
169x^4 - 39x^3 + 428x^2 + 89x + 203=0.
\]

You are told that, after relabeling if needed,
\[
r_1+r_2=\frac{8}{13}
\qquad\text{and}\qquad
r_3+r_4\ne r_1+r_2.
\]

Define the two \textbf{bond-prices}
\[
B_1:=r_1r_2,\qquad B_2:=r_3r_4.
\]
The Vault mints coins by the rule:

\begin{quote}
Write each bond-price \(B_i\) as a reduced fraction \(\frac{p_i}{q_i}\) with \(\gcd(p_i,q_i)=1\) and \(q_i>0\).
The bond costs \(p_i+q_i\) coins.
\end{quote}

\VaultDirective{Output the total number of coins required to buy both bonds.}

\SolutionPage
Let the elementary symmetric sums be
\[
s_1=r_1+r_2+r_3+r_4=-\frac{b}{a},\quad
s_2=\sum_{i<j}r_ir_j=\frac{c}{a},\quad
s_3=\sum_{i<j<k}r_ir_jr_k=-\frac{d}{a}.
\]
Given \(r_1+r_2=u=\frac{8}{13}\), and since \(s_1=-\frac{b}{a}=\frac{39}{169}=\frac{3}{13}\), we have \(r_3+r_4=v=s_1-u=-\frac{5}{13}\). In particular \(u\ne v\).

Set \(x=r_1r_2=B_1\), \(y=r_3r_4=B_2\).
Then
\[
s_2 = (r_1r_2+r_3r_4) + (r_1+r_2)(r_3+r_4)=x+y+uv,
\]
so \(x+y=s_2-uv\) is rational.
Also
\[
s_3=r_1r_2(r_3+r_4)+r_3r_4(r_1+r_2)=vx+uy,
\]
so \((x,y)\) solves the \(2\times2\) linear system
\[
\begin{cases}
x+y=s_2-uv,\\
vx+uy=s_3,
\end{cases}
\]
whose determinant \(u-v\neq0\). Hence \(x,y\in\mathbb Q\) uniquely.

Now plug in the given \(s_2=\frac{428}{169}\), \(s_3=-\frac{89}{169}\), \(u=\frac{8}{13}\), \(v=-\frac{5}{13}\):
\[
uv=-\frac{40}{169}\ \Rightarrow\ x+y=\frac{468}{169}.
\]
And \(vx+uy=s_3\) becomes \(-5x+8y=-\frac{89}{13}\).
Solving gives \(x=\frac{29}{13}\), \(y=\frac{7}{13}\).

Prices: \((29+13)+(7+13)=42+20=62\).

\FinalAnswer{62}

\NextPage
